% Projeto Integrador AOC 2026/02 --- artigo final, versão em português (formato IEEE)
% É a tradução de artigo.tex; os números vêm de numbers-pt.tex (vírgula decimal).
% Compilar: latexmk -pdf artigo-pt.tex     (ou, na raiz do repositório: make paper)
\documentclass[conference]{IEEEtran}

\usepackage[T1]{fontenc}
\usepackage[utf8]{inputenc}
\usepackage[brazilian]{babel}
\usepackage{amsmath,amssymb}
\usepackage{booktabs}
\usepackage{cite}
\usepackage{xcolor}
\usepackage{tikz}
\usetikzlibrary{arrows.meta,positioning,calc}
\usepackage{pgfplots}
\usepgfplotslibrary{groupplots}
\pgfplotsset{compat=1.18}
\usepackage[hidelinks]{hyperref}

\newcommand{\code}[1]{\texttt{#1}}
\newcommand{\res}{../../results}
\newcommand{\cdf}[3][]{\addplot+[no markers, #1] table[x=err, y=frac] {\res/cdf_#2_#3.dat};}

% Gerado por: python3 -m sim.paper_numbers (não editar à mão)
\newcommand{\ERRTABLE}{T8 & \boldmath$1{,}6\times10^{0}$ & $6{,}8\times10^{-2}$ & \boldmath$2{,}1\times10^{-1}$ & $7{,}7\times10^{-2}$ & $1{,}1\times10^{-1}$ & $1{,}4\times10^{-1}$ & $5{,}7\times10^{-2}$ & \boldmath$5{,}4\times10^{-1}$ & $1{,}4\times10^{-1}$\\
FP8 & $1{,}6\times10^{0}$ & $7{,}7\times10^{-2}$ & $\infty$ & $9{,}0\times10^{-2}$ & $5{,}4\times10^{-2}$ & $6{,}4\times10^{-2}$ & $4{,}4\times10^{-2}$ & $\infty$ & $7{,}2\times10^{-2}$\\
FP8w & $1{,}7\times10^{0}$ & \boldmath$4{,}4\times10^{-2}$ & $\infty$ & \boldmath$4{,}1\times10^{-2}$ & \boldmath$3{,}8\times10^{-2}$ & \boldmath$5{,}1\times10^{-2}$ & \boldmath$3{,}6\times10^{-2}$ & $\infty$ & \boldmath$6{,}4\times10^{-2}$\\
\midrule
T16 & \boldmath$6{,}9\times10^{-3}$ & \boldmath$2{,}7\times10^{-4}$ & $8{,}4\times10^{-4}$ & \boldmath$2{,}5\times10^{-4}$ & \boldmath$2{,}6\times10^{-4}$ & $5{,}4\times10^{-4}$ & \boldmath$2{,}3\times10^{-4}$ & \boldmath$2{,}0\times10^{-3}$ & $5{,}8\times10^{-4}$\\
FP16 & $9{,}1\times10^{-3}$ & $4{,}7\times10^{-4}$ & $3{,}8\times10^{-4}$ & $6{,}7\times10^{-4}$ & $3{,}9\times10^{-4}$ & $3{,}5\times10^{-4}$ & $2{,}9\times10^{-4}$ & $\infty$ & $3{,}8\times10^{-4}$\\
FP16w & $9{,}0\times10^{-3}$ & $3{,}2\times10^{-4}$ & \boldmath$2{,}4\times10^{-4}$ & $3{,}2\times10^{-4}$ & $3{,}0\times10^{-4}$ & \boldmath$3{,}1\times10^{-4}$ & $2{,}4\times10^{-4}$ & $\infty$ & \boldmath$3{,}1\times10^{-4}$\\
\midrule
T32 & \boldmath$9{,}9\times10^{-8}$ & \boldmath$4{,}2\times10^{-9}$ & \boldmath$1{,}4\times10^{-8}$ & \boldmath$3{,}4\times10^{-9}$ & \boldmath$3{,}7\times10^{-9}$ & \boldmath$8{,}3\times10^{-9}$ & \boldmath$3{,}8\times10^{-9}$ & \boldmath$3{,}3\times10^{-8}$ & \boldmath$8{,}6\times10^{-9}$\\
FP32 & $1{,}2\times10^{-6}$ & $6{,}9\times10^{-8}$ & $5{,}2\times10^{-8}$ & $1{,}0\times10^{-7}$ & $4{,}6\times10^{-8}$ & $5{,}2\times10^{-8}$ & $3{,}7\times10^{-8}$ & $6{,}2\times10^{-8}$ & $4{,}8\times10^{-8}$\\
FP32w & $9{,}3\times10^{-7}$ & $4{,}3\times10^{-8}$ & $3{,}1\times10^{-8}$ & $3{,}8\times10^{-8}$ & $3{,}5\times10^{-8}$ & $3{,}8\times10^{-8}$ & $3{,}0\times10^{-8}$ & $4{,}0\times10^{-8}$ & $3{,}8\times10^{-8}$\\}
\newcommand{\ERRTABLESYN}{T8 & \boldmath$1{,}6\times10^{0}$ & $6{,}8\times10^{-2}$ & \boldmath$2{,}1\times10^{-1}$ & $7{,}7\times10^{-2}$\\
FP8 & $1{,}6\times10^{0}$ & $7{,}7\times10^{-2}$ & $\infty$ & $9{,}0\times10^{-2}$\\
FP8w & $1{,}7\times10^{0}$ & \boldmath$4{,}4\times10^{-2}$ & $\infty$ & \boldmath$4{,}1\times10^{-2}$\\
\midrule
T16 & \boldmath$6{,}9\times10^{-3}$ & \boldmath$2{,}7\times10^{-4}$ & $8{,}4\times10^{-4}$ & \boldmath$2{,}5\times10^{-4}$\\
FP16 & $9{,}1\times10^{-3}$ & $4{,}7\times10^{-4}$ & $3{,}8\times10^{-4}$ & $6{,}7\times10^{-4}$\\
FP16w & $9{,}0\times10^{-3}$ & $3{,}2\times10^{-4}$ & \boldmath$2{,}4\times10^{-4}$ & $3{,}2\times10^{-4}$\\
\midrule
T32 & \boldmath$9{,}9\times10^{-8}$ & \boldmath$4{,}2\times10^{-9}$ & \boldmath$1{,}4\times10^{-8}$ & \boldmath$3{,}4\times10^{-9}$\\
FP32 & $1{,}2\times10^{-6}$ & $6{,}9\times10^{-8}$ & $5{,}2\times10^{-8}$ & $1{,}0\times10^{-7}$\\
FP32w & $9{,}3\times10^{-7}$ & $4{,}3\times10^{-8}$ & $3{,}1\times10^{-8}$ & $3{,}8\times10^{-8}$\\}
\newcommand{\ERRTABLEREAL}{T8 & $1{,}1\times10^{-1}$ & $1{,}4\times10^{-1}$ & $5{,}7\times10^{-2}$ & \boldmath$5{,}4\times10^{-1}$ & $1{,}4\times10^{-1}$\\
FP8 & $5{,}4\times10^{-2}$ & $6{,}4\times10^{-2}$ & $4{,}4\times10^{-2}$ & $\infty$ & $7{,}2\times10^{-2}$\\
FP8w & \boldmath$3{,}8\times10^{-2}$ & \boldmath$5{,}1\times10^{-2}$ & \boldmath$3{,}6\times10^{-2}$ & $\infty$ & \boldmath$6{,}4\times10^{-2}$\\
\midrule
T16 & \boldmath$2{,}6\times10^{-4}$ & $5{,}4\times10^{-4}$ & \boldmath$2{,}3\times10^{-4}$ & \boldmath$2{,}0\times10^{-3}$ & $5{,}8\times10^{-4}$\\
FP16 & $3{,}9\times10^{-4}$ & $3{,}5\times10^{-4}$ & $2{,}9\times10^{-4}$ & $\infty$ & $3{,}8\times10^{-4}$\\
FP16w & $3{,}0\times10^{-4}$ & \boldmath$3{,}1\times10^{-4}$ & $2{,}4\times10^{-4}$ & $\infty$ & \boldmath$3{,}1\times10^{-4}$\\
\midrule
T32 & \boldmath$3{,}7\times10^{-9}$ & \boldmath$8{,}3\times10^{-9}$ & \boldmath$3{,}8\times10^{-9}$ & \boldmath$3{,}3\times10^{-8}$ & \boldmath$8{,}6\times10^{-9}$\\
FP32 & $4{,}6\times10^{-8}$ & $5{,}2\times10^{-8}$ & $3{,}7\times10^{-8}$ & $6{,}2\times10^{-8}$ & $4{,}8\times10^{-8}$\\
FP32w & $3{,}5\times10^{-8}$ & $3{,}8\times10^{-8}$ & $3{,}0\times10^{-8}$ & $4{,}0\times10^{-8}$ & $3{,}8\times10^{-8}$\\}
\newcommand{\RATA}{1{,}2}
\newcommand{\RATB}{11}
\newcommand{\WIDET}{$8{,}4\times10^{-4}$}
\newcommand{\WIDEF}{$2{,}4\times10^{-4}$}
\newcommand{\BUST}{$5{,}4\times10^{-4}$}
\newcommand{\BUSF}{$3{,}1\times10^{-4}$}
\newcommand{\WIDETEIGHT}{$2{,}1\times10^{-1}$}
\newcommand{\BCSTEIGHT}{$5{,}4\times10^{-1}$}
\newcommand{\BCSTSIXTEEN}{$2{,}0\times10^{-3}$}
\newcommand{\FPEIGHTBAD}{217}
\newcommand{\NFBUS}{94}
\newcommand{\NFWEST}{51}
\newcommand{\NFBCS}{153}
\newcommand{\NFWESTH}{5}
\newcommand{\NFBCSH}{149}
\newcommand{\MEMTE}{$1{,}4\times10^{-1}$}
\newcommand{\MEMTEBAD}{267}
\newcommand{\MEMFE}{$6{,}4\times10^{-2}$}
\newcommand{\MEMFEBAD}{1\,054}
\newcommand{\MEMTS}{$5{,}8\times10^{-4}$}
\newcommand{\MEMTSBAD}{3}
\newcommand{\MEMFS}{$3{,}1\times10^{-4}$}
\newcommand{\MEMFSBAD}{2}
\newcommand{\MEMTT}{$8{,}6\times10^{-9}$}
\newcommand{\MEMTTBAD}{0}
\newcommand{\MEMFT}{$3{,}8\times10^{-8}$}
\newcommand{\MEMFTBAD}{0}
\newcommand{\MEMROWS}{17\,758}
\newcommand{\MATTABLE}{poisson & 256$\times$256 & 1\,216 & 4{,}8 & $2^{0}$ -- $2^{2}$\\
gauss & 256$\times$256 & 2\,560 & 10{,}0 & $2^{-13}$ -- $2^{2}$\\
wide & 256$\times$256 & 2\,560 & 10{,}0 & $2^{-10}$ -- $2^{10}$\\
mlp & 256$\times$512 & 13\,056 & 51{,}0 & $2^{-15}$ -- $2^{0}$\\
\midrule
nos4 & 100$\times$100 & 594 & 5{,}9 & $2^{-6}$ -- $2^{-1}$\\
bus494 & 494$\times$494 & 1\,666 & 3{,}4 & $2^{-3}$ -- $2^{14}$\\
west0479 & 479$\times$479 & 1\,910 & 4{,}0 & $2^{-21}$ -- $2^{18}$\\
bcsstk05 & 153$\times$153 & 2\,423 & 15{,}8 & $2^{-31}$ -- $2^{22}$\\
memplus & 17758$\times$17758 & 126\,150 & 7{,}1 & $2^{-86}$ -- $2^{0}$\\}
\newcommand{\INSTRTABLE}{poisson & 4{,}8 & 6\,146 & 6\,146 & 6\,658 & 0{,}0\% & 7{,}7\%\\
gauss & 10{,}0 & 8\,962 & 9\,218 & 9\,986 & 2{,}8\% & 10{,}3\%\\
mlp & 51{,}0 & 23\,042 & 24\,578 & 26\,626 & 6{,}2\% & 13{,}5\%\\
\midrule
nos4 & 5{,}9 & 2\,402 & 2\,402 & 2\,602 & 0{,}0\% & 7{,}7\%\\
bus494 & 3{,}4 & 11\,902 & 11\,906 & 12\,898 & 0{,}0\% & 7{,}7\%\\
west0479 & 4{,}0 & 11\,949 & 11\,990 & 12\,989 & 0{,}3\% & 8{,}0\%\\
bcsstk05 & 15{,}8 & 6\,061 & 6\,278 & 6\,801 & 3{,}5\% & 10{,}9\%\\
memplus & 7{,}1 & 503\,161 & 510\,158 & 552\,671 & 1{,}4\% & 9{,}0\%\\}
\newcommand{\SAVEWLO}{7{,}7}
\newcommand{\SAVEWHI}{14{,}3}
\newcommand{\SAVEPHI}{7{,}1}
\newcommand{\FLAGPCT}{90}
\newcommand{\FLAGWHERE}{\emph{wide}, FP32}
\newcommand{\HWDECA}{46}
\newcommand{\HWDECB}{89}
\newcommand{\HWDECC}{164}
\newcommand{\HWENCA}{70}
\newcommand{\HWENCB}{118}
\newcommand{\HWENCC}{216}
\newcommand{\HWALUA}{327}
\newcommand{\HWALUB}{912}
\newcommand{\HWALUC}{2\,838}
\newcommand{\HWNCA}{32}
\newcommand{\HWNCB}{67}
\newcommand{\HWNCC}{155}
\newcommand{\HWWMA}{220}
\newcommand{\HWWMB}{686}
\newcommand{\HWWMC}{2\,384}
\newcommand{\HWDOTA}{6\,740}
\newcommand{\HWDOTB}{7\,525}
\newcommand{\HWDOTC}{9\,451}
\newcommand{\HWHDR}{41}

\newcommand{\matplot}[2][]{%
  \nextgroupplot[title={#2}, #1]
  \cdf[blue!75!black]{#2}{T8}    \cdf[blue!75!black, dashed]{#2}{FP8w}
  \cdf[red!75!black]{#2}{T16}    \cdf[red!75!black, dashed]{#2}{FP16w}
  \cdf[green!45!black]{#2}{T32}  \cdf[green!45!black, dashed]{#2}{FP32w}}

\renewcommand{\IEEEkeywordsname}{Palavras-chave}

\begin{document}

\title{Zvtakum: Aritmética Takum na Extensão Vetorial RISC-V\\
{\large Caminho de Dados em RTL, Simulação em Nível de Instrução e Avaliação em Álgebra Linear Esparsa}}

\author{\IEEEauthorblockN{Heloisa Silva Rosa, Samantha Kelly de Souza Sena e Vitor Hugo Lacerda Lana}
\IEEEauthorblockA{Departamento de Engenharia Eletrônica (DELT), Universidade Federal de Minas Gerais\\
Belo Horizonte, MG, Brasil}}

\maketitle

\begin{abstract}
Os conjuntos de instruções vetoriais continuam acrescentando formatos de baixa precisão no estilo
IEEE~754, cada um com sua lógica de decodificação, seus casos de exceção e seu estado de controle. A
aritmética takum é uma alternativa de precisão afunilada cujo cabeçalho é decodificado a partir de, no
máximo, 12 bits em qualquer largura. Apresentamos a Zvtakum, uma extensão takum simplificada para a
extensão vetorial do RISC-V, com cinco instruções (\code{vtadd.vv}, \code{vtsub.vv}, \code{vtmul.vv},
\code{vtdot.vv}, \code{vncvt.t.t}) para elementos de 8, 16 e 32 bits. Descrevemos um caminho de dados em
Verilog, verificado contra um oráculo racional exato (exaustivamente em 8 bits), e um simulador em nível
de instrução que executa as instruções codificadas junto com um subconjunto do RVV. No produto de matriz
esparsa por vetor com nove matrizes, cinco delas do repositório público Matrix Market e uma com 17\,758
linhas, comparado com uma referência
em binary128, os takums de 32 bits têm o menor erro relativo mediano em todos os casos, até 11 vezes menor
que o do binary32 com a mesma precisão de acúmulo; os de 16 bits ganham com dados próximos da escala
unitária, perdem em faixas moderadamente largas e permanecem finitos onde o binary16 estoura; e os de 8
bits são menos precisos que um ponto flutuante de 8 bits com dados dentro da faixa, mas nunca estouram,
enquanto o ponto flutuante devolve infinitos ou NaN em 11\% a 100\% das linhas de quatro matrizes. A instrução de produto escalar fundido
elimina de 8 a 14\% das instruções executadas em relação ao núcleo IEEE equivalente com alargamento, e a
extensão dispensa o registrador de controle e status de ponto flutuante.
\end{abstract}

\begin{IEEEkeywords}
aritmética takum, RISC-V, extensão vetorial, aritmética computacional, álgebra linear esparsa
\end{IEEEkeywords}

%======================================================================
\section{Introdução}

As unidades SIMD (\emph{single instruction, multiple data}) e vetoriais foram projetadas em torno do
binary32 e do binary64. As cargas de aprendizado de máquina as empurraram para formatos mais estreitos:
binary16, bfloat16 e, mais recentemente, os formatos de 8 bits E4M3 e E5M2~\cite{ofp8}. Cada um tem uma
largura de expoente diferente, um tratamento diferente de infinitos e NaN e, às vezes, uma regra de
saturação diferente; um conjunto de instruções que suporte todos acumula instruções de conversão, casos
especiais e lógica de exceção~\cite{hunhold2025avx}.

A aritmética takum~\cite{hunhold2024takum} substitui essa família por um único formato de precisão
afunilada, definido para qualquer largura $n$. Os campos de sinal, direção, regime e característica cabem
sempre nos 12 bits mais significativos, a faixa dinâmica é de cerca de $2^{\pm 239}$ já em 8 bits e de $2^{\pm 255}$ a partir de 12, há uma
única representação do zero e um único valor excepcional (NaR), e os padrões de bits são ordenados como
inteiros em complemento de dois.

Este trabalho avalia o que isso significa para um conjunto de instruções vetoriais. As contribuições são:
\begin{itemize}
  \item um caminho de dados combinacional em Verilog para takums lineares (decodificador, codificador,
        ALU de soma, subtração e multiplicação, multiplicador com alargamento, conversão com redução de
        largura e uma unidade vetorial), parametrizado para T8, T16 e T32;
  \item a Zvtakum, uma codificação e uma semântica para cinco instruções vetoriais takum que espelham
        instruções de ponto flutuante do RVV, incluindo um produto escalar fundido com acúmulo em
        precisão ampliada;
  \item um simulador em nível de instrução de RV64I, de um subconjunto do RVV~1.0 e da Zvtakum,
        conferido contra os vetores de teste do RTL;
  \item uma comparação experimental com formatos no estilo IEEE~754 no produto de matriz esparsa por
        vetor: erro relativo contra binary128, instruções executadas e requisitos de estado de controle.
\end{itemize}
Todo o código, os testes e os dados estão no repositório do projeto e são reproduzidos por
\code{make test} e \code{make bench}.

%======================================================================
\section{Trabalhos Relacionados}

Os posits~\cite{gustafson2017} introduziram a precisão afunilada com um regime de comprimento variável e
foram depois padronizados~\cite{positstd}. O regime de um posit pode consumir a palavra inteira, o que faz
a faixa dinâmica depender de $n$ e o decodificador depender de uma contagem de bits à esquerda na largura
total. Os takums~\cite{hunhold2024takum} limitam o regime a três bits e a característica a sete, de modo
que o expoente é sempre recuperado de um cabeçalho fixo de 12 bits; um codec em VHDL construído sobre
essa propriedade é descrito em~\cite{hunhold2024codec}.

Hunhold~\cite{hunhold2025avx} estudou o Intel AVX10.2 e mostrou que substituir seus formatos derivados do
IEEE por takums remove ou generaliza grande parte do conjunto de instruções, pois comparações, manipulação
de sinal e conversões entre larguras se reduzem a operações inteiras. Nosso trabalho faz a mesma pergunta
para a extensão vetorial do RISC-V~\cite{rvv}, em escala menor, e acrescenta uma implementação. No lado
numérico, os takums já foram comparados com bfloat16 e posits em solucionadores lineares
esparsos~\cite{hunhold2025solvers}; usamos o mesmo tipo de experimento (erro contra uma referência de alta
precisão), restrito ao produto matriz--vetor, com matrizes sintéticas e com matrizes do
repositório Matrix Market~\cite{boisvert1997}, a maioria da coleção Harwell-Boeing~\cite{duff1989}.

%======================================================================
\section{Formato Takum}

Usamos o takum linear~\cite{hunhold2024takum}. Uma palavra de $n$ bits contém um sinal $S$, um bit de
direção $D$, um regime $R$ de 3 bits, uma característica $C$ de $r$ bits e uma fração $F$ de $p=n-5-r$
bits, com
\begin{equation}
r=\begin{cases}R & D=1\\ 7-R & D=0\end{cases}\qquad
c=\begin{cases}2^{r}-1+C & D=1\\ -2^{r+1}+1+C & D=0\end{cases}
\end{equation}
e, para uma palavra positiva, $x=(1+f)\,2^{c}$ com $f=2^{-p}F$. Os valores negativos são o complemento de
dois da palavra. O padrão $0\cdots0$ é o zero e $10\cdots0$ é o NaR. A característica $c$ vai de $-255$ a
$254$. A precisão é afunilada: valores próximos de 1 têm $r=0$ e $n-5$ bits de fração; as maiores
magnitudes têm $r=7$ e $n-12$. A Tabela~\ref{tab:formats} compara os bits de fração com os dos formatos
no estilo IEEE usados neste artigo.

\begin{table}[t]
\caption{Bits de fração e maior expoente finito de cada formato.}
\label{tab:formats}
\centering
\setlength{\tabcolsep}{4pt}
\begin{tabular}{@{}lccc@{}}
\toprule
Formato & $n$ & Bits de fração & Expoente máx.\\
\midrule
T8, T16, T32 & 8, 16, 32 & $n-5-r$,\; $0\le r\le 7$ & 239, 254, 254\\
FP8 (E4M3) & 8 & 3 & 7\\
FP16 (binary16) & 16 & 10 & 15\\
FP32 (binary32) & 32 & 23 & 127\\
\bottomrule
\end{tabular}
\end{table}

%======================================================================
\section{Arquitetura Proposta}

\subsection{Caminho de dados escalar}

Toda operação é \emph{decodificar $\rightarrow$ núcleo $\rightarrow$ codificar}. O decodificador nega
condicionalmente a palavra, lê o cabeçalho de 12 bits com um bloco fixo (\code{takum\_hdr\_dec},
compartilhado por todas as larguras) e alinha a fração à esquerda com um único deslocador. O núcleo opera
sobre sinal, um expoente de 9 bits e um significando: a multiplicação soma os expoentes e multiplica os
significandos; a soma ordena os operandos por magnitude, alinha o menor preservando os bits de guarda,
arredondamento e \emph{sticky}, soma ou subtrai e normaliza. O codificador reconstrói o cabeçalho, anexa a
fração e arredonda \emph{na palavra}: como os takums são monótonos no padrão de bits, somar um à palavra
truncada leva ao próximo valor representável, inclusive ao cruzar fronteiras de regime. O arredondamento é
para o mais próximo, com empate para o par, e um resultado não nulo satura em vez de virar zero ou NaR.

\subsection{Instruções da Zvtakum}

A Tabela~\ref{tab:isa} lista as instruções. Elas usam o layout de campos do OP-V
(\code{funct6 | vm | vs2 | vs1 | funct3 | vd | opcode}) no opcode principal \emph{custom-2}, com o
\code{funct6} da instrução de ponto flutuante que substituem, de modo que o decodificador de um núcleo
RVV pode ser compartilhado. A largura do elemento, o comprimento do vetor e a máscara vêm dos já
existentes \code{vtype}, \code{vl} e \code{v0}; a extensão não acrescenta estado.

\begin{table}[t]
\caption{Instruções da Zvtakum. SEW é 8, 16 ou 32.}
\label{tab:isa}
\centering
\begin{tabular}{@{}lll@{}}
\toprule
Instrução & \code{funct6} & Operação ($i<\mathit{vl}$, não mascarado)\\
\midrule
\code{vtadd.vv} & \code{000000} & $vd_i = vs2_i + vs1_i$\\
\code{vtsub.vv} & \code{000010} & $vd_i = vs2_i - vs1_i$\\
\code{vtmul.vv} & \code{100100} & $vd_i = vs2_i \times vs1_i$\\
\code{vtdot.vv} & \code{111001} & $vd_0^{(2\mathrm{SEW})} \mathrel{+}= \sum_i vs2_i \times vs1_i$\\
\code{vncvt.t.t} & \code{010010} & $vd_i = \mathrm{arred}_{\mathrm{SEW}}\bigl(vs2_i^{(2\mathrm{SEW})}\bigr)$\\
\bottomrule
\end{tabular}
\end{table}

\paragraph*{Redução de largura} Nos formatos IEEE, estreitar exige refazer o viés do expoente, tratar
subnormais e detectar estouro. Um takum de $2n$ bits é estreitado arredondando sua cadeia de bits para $n$
bits: o cabeçalho fica no topo da palavra em qualquer largura. O \code{takum\_ncvt} é um incrementador e
um teste de saturação, sem decodificador.

\paragraph*{Produto escalar com acúmulo ampliado} O \code{vtdot.vv} (Fig.~\ref{fig:dot}) multiplica os
elementos de SEW bits em takums de $2\,$SEW bits e os acumula, na ordem dos elementos, no elemento 0 de
\code{vd}, também em $2\,$SEW. O produto de dois significandos de $n-4$ bits tem $2n-8$ bits e sempre cabe
na fração de um takum de $2n$ bits, de modo que o multiplicador com alargamento (\code{takum\_wmul}) é
exato, salvo quando o expoente satura; conferimos isso para todos os $2^{16}$ pares de T8. O
arredondamento ocorre apenas nas somas, em largura dupla, e uma vez mais no \code{vncvt.t.t} final. Como o
acumulador é lido de \code{vd}, um laço com \code{vsetvli} não precisa de redução separada nem de
movimentação escalar.

\begin{figure}[t]
\centering
\begin{tikzpicture}[x=0.92cm, font=\footnotesize, >={Latex[length=1.6mm]},
  blk/.style={draw, rounded corners=1.5pt, minimum height=6.5mm, minimum width=13mm, align=center}]
  \foreach \i/\x in {0/0, 1/2.1, 2/4.9} {
    \node[blk] (m\i) at (\x, 0) {wmul};
    \node[blk] (a\i) at (\x,-1.25) {soma$_{2n}$};
    \draw[->] (m\i) -- (a\i);
  }
  \node at (3.5, 0) {$\cdots$}; \node at (3.5,-1.25) {$\cdots$};
  \node (in0) at (0, 0.95) {$a_0\; b_0$}; \node (in1) at (2.1, 0.95) {$a_1\; b_1$};
  \node (in2) at (4.9, 0.95) {$a_{k-1}\; b_{k-1}$};
  \foreach \i in {0,1,2} \draw[->] (in\i) -- (m\i);
  \node (acc) at (-1.5,-1.25) {$vd_0$};
  \draw[->] (acc) -- (a0); \draw[->] (a0) -- (a1);
  \draw[->] (a1) -- (2.95,-1.25); \draw[->] (4.05,-1.25) -- (a2);
  \node[blk] (nc) at (6.75,-1.25) {ncvt};
  \draw[->] (a2) -- node[above] {$2n$} (nc);
  \node (y) at (6.75,-2.2) {$n$ bits};
  \draw[->] (nc) -- (y);
  \node[align=center, text=black!60] at (2.45,-2.05) {\scriptsize a pista é contornada se $i\ge \mathit{vl}$ ou mascarada};
\end{tikzpicture}
\caption{Caminho de dados do \code{vtdot.vv} para $k$ pistas de elementos de $n$ bits, seguido do
\code{vncvt.t.t}. Os produtos são exatos em $2n$ bits; só as somas arredondam.}
\label{fig:dot}
\end{figure}

\subsection{Unidade vetorial e simulador}

A unidade vetorial \code{takum\_vec\_alu} (parâmetro VLEN, 64 por padrão) instancia uma fileira de ALUs
por largura de elemento e seleciona pela SEW. Os elementos com $i\ge \mathit{vl}$ e os mascarados mantêm o
valor antigo do destino. Ela é combinacional; uma implementação com \emph{pipeline} registraria entre os
três estágios e compartilharia os multiplicadores do \code{vtmul.vv} e do \code{vtdot.vv}.

Para executar programas inteiros, escrevemos um simulador em nível de instrução em Python. Ele monta um
pequeno subconjunto de RV64I, as instruções do RVV de configuração, carga e armazenamento, carga indexada
e ponto flutuante usadas pelos núcleos, e a Zvtakum; a máquina busca palavras de 32 bits, decodifica,
executa e conta as instruções executadas. Os formatos IEEE~754~\cite{ieee754} são modelados por um formato
binário genérico corretamente arredondado (subnormais, infinitos, NaN e os cinco \code{fflags}); o zero
negativo não é modelado. As operações takum chamam o mesmo oráculo que verifica o RTL.

%======================================================================
\section{Resultados e Discussão}

\subsection{Verificação}

O RTL é comparado com um oráculo que decodifica os takums pela fórmula completa e calcula em aritmética
racional exata, sem compartilhar código nem método com o hardware. A Tabela~\ref{tab:verif} resume a suíte
de testes; todos passam. Para conferir que os testes são sensíveis, injetamos 13 defeitos (por exemplo,
descartar o bit \emph{sticky}, omitir a regra de empate para o par na redução de largura, acumular o
produto escalar em ordem inversa, ignorar a máscara); todos foram detectados. O oráculo, por sua vez, foi comparado com a libtakum~\cite{libtakum}, a
implementação de referência do autor do formato: 826\,987 vetores coincidem bit a bit, e 16 produtos em T32
diferem no último bit porque a libtakum multiplica em binary64, que não comporta o produto de 56 bits; nos
16, o resultado exato está mais perto do nosso. O simulador é conferido
contra a aritmética nativa em binary16/32/64 (160\,000 casos) e executa os 15\,000 vetores de teste da
unidade vetorial como instruções codificadas, com resultados idênticos.

\begin{table}[t]
\caption{Verificação do RTL contra o oráculo exato.}
\label{tab:verif}
\centering
\begin{tabular}{@{}lr@{}}
\toprule
Teste & Vetores\\
\midrule
Decodificador, codificador, ida e volta (T8/T16 exaustivos) & 855\,408\\
ALU T8: todos os pares $\times$ soma, subtração, produto & 196\,608\\
ALU T16, T32: bordas, vizinhos, aleatórios (cada) & 109\,923\\
Redução T16$\to$T8 (todas as palavras), T32$\to$T16, T64$\to$T32 & 105\,794\\
Produto alargado T8 (todos os pares), T16, T32 & 138\,818\\
Unidade vetorial, 5 instruções, \code{vl} e máscara (por SEW) & 5\,000\\
\bottomrule
\end{tabular}
\end{table}

\subsection{Custo de hardware}

A Tabela~\ref{tab:hw} traz uma síntese independente de tecnologia (Yosys, mapeamento em LUTs de 6
entradas). Os números servem para comparar os blocos entre si, não como área absoluta. O decodificador do
cabeçalho custa \HWHDR{} LUTs em qualquer largura. A conversão com redução de largura é menor que um
decodificador, o que confirma que ela é apenas um arredondamento da palavra. O multiplicador com
alargamento é mais barato que a ALU completa porque não tem deslocador de alinhamento nem de normalização;
em 32 bits o multiplicador $28\times28$ domina ambos. O \code{vtdot.vv} combinacional é o maior bloco,
pois cada pista carrega um somador de largura dupla; um projeto preocupado com área iteraria um único
somador sobre as pistas.

\begin{table}[t]
\caption{Resultados de síntese em LUT6 (Yosys, genérica).}
\label{tab:hw}
\centering
\begin{tabular}{lrrr}
\toprule
Bloco & T8 & T16 & T32\\
\midrule
Decodificador & \HWDECA & \HWDECB & \HWDECC\\
Codificador & \HWENCA & \HWENCB & \HWENCC\\
ALU (soma, subtração, produto) & \HWALUA & \HWALUB & \HWALUC\\
Redução para esta largura (\code{vncvt.t.t}) & \HWNCA & \HWNCB & \HWNCC\\
Multiplicador com alargamento & \HWWMA & \HWWMB & \HWWMC\\
\code{vtdot.vv}, vetor de 64 bits (8/4/2 pistas) & \HWDOTA & \HWDOTB & \HWDOTC\\
\bottomrule
\end{tabular}
\end{table}

\subsection{Exatidão do produto de matriz esparsa por vetor}

Calculamos $y=Ax$ para as nove matrizes da Tabela~\ref{tab:mats}, em formato CSR. Quatro são geradas:
\emph{poisson} é o Laplaciano de 5 pontos em uma malha $16\times16$ aplicado a um modo suave, com forte
cancelamento em todas as linhas; \emph{gauss} é esparsa aleatória com entradas $\mathcal{N}(0,1)$;
\emph{wide} tem a mesma estrutura, com magnitudes log-uniformes; e \emph{mlp} é uma camada densa podada
(90\% de zeros, pesos $\mathcal{N}(0,2/\mathrm{fan\_in})$) aplicada a ativações pós-ReLU, um núcleo típico
de aprendizado de máquina. Quatro são problemas reais da coleção Harwell-Boeing~\cite{duff1989}, obtidos
do Matrix Market~\cite{boisvert1997}: \emph{nos4} (análise estrutural de uma viga), \emph{bus494}
(\code{494\_bus}, matriz de admitância de uma rede elétrica de 494 barras), \emph{west0479} (engenharia
química, uma coluna de destilação) e \emph{bcsstk05} (matriz de rigidez estrutural). Uma quinta matriz real
do mesmo repositório, \emph{memplus} (simulação de um circuito de memória de computador), tem 17\,758
linhas, 36 vezes a maior das demais, e entradas tão pequenas quanto $2^{-86}$. Para as
matrizes reais, $x$ é sorteado de $\mathcal{N}(0,1)$. Os dados ficam em binary64, são arredondados para cada
formato e multiplicados pelo núcleo em assembly no simulador. A referência é o mesmo produto calculado em
binary128 a partir dos dados não arredondados, e a métrica é o erro relativo de cada linha.

\begin{table}[t]
\caption{Matrizes de teste. Acima: geradas; abaixo: Matrix Market.}
\label{tab:mats}
\centering
\setlength{\tabcolsep}{4pt}
\begin{tabular}{@{}lcrrc@{}}
\toprule
Matriz & Dimensão & Não nulos & Por linha & Faixa de $|a_{ij}|$\\
\midrule
\MATTABLE
\bottomrule
\end{tabular}
\end{table}

Três núcleos são comparados: \emph{T}$n$ usa \code{vtdot.vv} e \code{vncvt.t.t}; \emph{FP}$n$ usa
\code{vfmul.vv} e \code{vfredosum.vs} em SEW, a formulação usual em RVV; e \emph{FP}$n$\emph{w} usa
\code{vfwmul.vv}, uma soma ordenada em $2\,$SEW e \code{vfncvt.f.f.w}, o que dá ao IEEE o mesmo acúmulo
ampliado do \code{vtdot.vv}. A Tabela~\ref{tab:err} traz o erro mediano e a Fig.~\ref{fig:cdf}, as
distribuições acumuladas.

\begin{table*}[t]
\caption{Erro relativo mediano de $y=Ax$ contra binary128. Negrito: o melhor de cada largura. $\infty$: mais
da metade das linhas não é finita.}
\label{tab:err}
\centering
\setlength{\tabcolsep}{4.5pt}
\begin{tabular}{@{}lccccccccc@{}}
\toprule
Núcleo & poisson & gauss & wide & mlp & nos4 & bus494 & west0479 & bcsstk05 & memplus\\
\midrule
\ERRTABLE
\bottomrule
\end{tabular}
\end{table*}

\begin{figure*}[t]
\centering
\begin{tikzpicture}
\begin{groupplot}[
  group style={group size=5 by 2, horizontal sep=4mm, vertical sep=11mm, ylabels at=edge left,
               yticklabels at=edge left, xlabels at=edge bottom},
  width=0.232\textwidth, height=36mm,
  xmode=log, xmin=1e-10, xmax=1e2, ymin=0, ymax=1,
  xtick={1e-10,1e-8,1e-6,1e-4,1e-2,1e0,1e2},
  xticklabels={,$10^{-8}$,,$10^{-4}$,,$10^{0}$,},
  yticklabels={,0,{0,2},{0,4},{0,6},{0,8},1},
  xlabel={erro relativo}, ylabel={fração das linhas},
  grid=major, grid style={black!12},
  tick label style={font=\scriptsize}, label style={font=\footnotesize}, title style={font=\footnotesize, yshift=-1.5mm},
  cycle list={}, every axis plot/.append style={line width=0.8pt},
  legend style={font=\scriptsize, legend columns=6, /tikz/every even column/.append style={column sep=3mm}}]
\matplot{poisson}
\matplot{gauss}
\matplot{wide}
\matplot{mlp}
\matplot{nos4}
\matplot{bus494}
\matplot{west0479}
\matplot{bcsstk05}
\matplot[legend to name=cdfleg]{memplus}
\legend{T8, FP8w, T16, FP16w, T32, FP32w}
\end{groupplot}
\end{tikzpicture}\\[1.5mm]
\pgfplotslegendfromname{cdfleg}
\caption{Distribuição acumulada do erro relativo de $y=Ax$ sobre as linhas de cada matriz, contra
binary128. Contínuo: takum com \code{vtdot.vv}; tracejado: IEEE com acúmulo ampliado. Curvas mais à
esquerda são melhores. Uma curva que não chega a 1 tem resultados não finitos.}
\label{fig:cdf}
\end{figure*}

\emph{32 bits.} O T32 tem o menor erro mediano nas nove matrizes, por um fator de \RATA{} a \RATB{} em
relação ao FP32w. A maioria dos operandos tem $r\le3$, de modo que o T32 carrega de 24 a 27 bits de
fração contra os 23 do binary32; o menor ganho ocorre em \emph{bcsstk05}, cujas entradas chegam a
$2^{22}$ ($r=4$, 23 bits).

\emph{16 bits.} O T16 é melhor que o FP16w em \emph{poisson}, \emph{gauss}, \emph{mlp}, \emph{nos4} e
\emph{west0479}, em que a maior parte dos dados está a poucas binadas de 1 e o T16 tem 10 ou 11 bits de
fração. Ele é pior em \emph{wide} (\WIDET{} contra \WIDEF{}), em \emph{bus494} (\BUST{} contra
\BUSF{}) e em \emph{memplus} (\MEMTS{} contra \MEMFS{}): magnitudes em torno de $2^{\pm10}$ exigem $r=3$, restando 8 bits de fração contra os 10 do
binary16. Esse é o custo do afunilamento, e ele só aparece enquanto o binary16 ainda tem faixa. Além dela
a comparação se inverte: o binary16 devolve infinitos ou NaN em \NFWESTH{} das 479 linhas de
\emph{west0479} e em \NFBCSH{} das 153 linhas de \emph{bcsstk05}, em que o T16 é finito em todas, com
erro mediano de \BCSTSIXTEEN.

\emph{8 bits.} Onde o ponto flutuante de 8 bits permanece na faixa, ele é o formato mais preciso, por um
fator de 1,5 a 3 na mediana: o T8 tem três bits de fração apenas em $[0{,}5;\,2)$, dois nas binadas
vizinhas e menos além delas. Ambos são grosseiros demais para \emph{poisson}, em que o cancelamento
amplifica o arredondamento das entradas para erros acima de 100\% na maioria das linhas. A vantagem do T8
é a faixa. O FP8w devolve infinitos ou NaN em \FPEIGHTBAD{} das 256 linhas de \emph{wide}, \NFBUS{} das
494 de \emph{bus494}, \NFWEST{} das 479 de \emph{west0479} e em todas as \NFBCS{} de \emph{bcsstk05}, e
nenhum núcleo evita isso, porque os próprios dados estão fora da faixa. Em \emph{memplus} o limite está na
outra ponta: entradas de até $2^{-86}$ causam \emph{underflow} no ponto flutuante de 8 bits, cujo resultado
erra 100\% ou mais em \MEMFEBAD{} das \MEMROWS{} linhas, contra \MEMTEBAD{} do T8, embora o erro mediano
do ponto flutuante seja menor. O T8 devolve um resultado finito em todas as linhas de todas as matrizes.

\emph{Efeito do acúmulo ampliado.} Repetindo o produto takum com \code{vtmul.vv} e \code{vtadd.vv} em SEW,
o erro mediano cresce por um fator de 1,2 a 2,3. O estreitamento final do \code{vtdot.vv} é um único
arredondamento, de modo que o erro restante é dominado pelo arredondamento das entradas: nas oito matrizes
sem cancelamento sistemático, o melhor resultado possível (a referência corretamente arredondada) é menos
de 3 vezes mais exato.

\subsection{Contagem de instruções}

\begin{figure}[t]
\footnotesize
\begin{verbatim}
inner:
  vsetvli    t3, t2, e16, m1, ta, ma
  vle32.v    v8, (a2)        # colunas
  vle16.v    v1, (a3)        # valores da linha
  vluxei32.v v2, (a4), v8    # x[col] (indexada)
  vtdot.vv   v4, v1, v2      # T32 += v1 . v2
  ...                        # 6 instr. de laco
done:
  vsetvli    zero, s1, e16, m1, ta, ma  # vl=1
  vncvt.t.t  v5, v4          # T32 -> T16
  vse16.v    v5, (a5)

# IEEE com o mesmo acumulo ampliado,
# no lugar do vtdot.vv:
  vfwmul.vv    v6, v1, v2    # produtos em binary32
  vsetvli      zero, t3, e32, m2, ta, ma
  vfredosum.vs v4, v6, v4
\end{verbatim}
\caption{Laço interno e fecho da linha do núcleo CSR para elementos de 16 bits.}
\label{fig:asm}
\end{figure}

A Fig.~\ref{fig:asm} mostra o núcleo e a Tabela~\ref{tab:instr} conta as instruções executadas. O laço
interno tem 11 instruções por faixa com \code{vtdot.vv}, 12 com \code{vfmul}/\code{vfredosum} e 13 com
alargamento, que precisa de um segundo \code{vsetvli} por faixa porque os produtos e a redução usam
larguras de elemento diferentes. O \code{vtdot.vv} produz a soma ampliada sem alterar o \code{vtype}.
Considerando todas as matrizes e larguras, a economia frente ao núcleo IEEE da mesma classe de exatidão
(FPw) vai de \SAVEWLO{} a \SAVEWHI\% com VLEN${}=128$; frente ao núcleo simples, de 0 a \SAVEPHI\%. A
economia cresce com o número de não nulos por linha. Quando a linha cabe em uma única faixa, domina o
custo fixo por linha (12 ou 13 instruções) e o núcleo takum executa tantas instruções quanto o núcleo
IEEE simples, com a instrução fundida pagando a conversão final; as matrizes reais, exceto \emph{bcsstk05},
têm de 3 a 7 não nulos por linha e ficam nesse caso ou perto dele.

\begin{table}[t]
\caption{Instruções executadas para $y=Ax$, SEW${}=16$, VLEN${}=128$. \emph{wide} tem a estrutura de
\emph{gauss}.}
\label{tab:instr}
\centering
\setlength{\tabcolsep}{3.2pt}
\begin{tabular}{@{}lrrrrrr@{}}
\toprule
 & nnz/ & & & & \multicolumn{2}{c}{Economia vs.}\\
Matriz & linha & Zvtakum & FP & FPw & FP & FPw\\
\midrule
\INSTRTABLE
\bottomrule
\end{tabular}
\end{table}

\subsection{Complexidade do conjunto de instruções e dos CSRs}

Além da contagem de instruções, a unidade takum elimina estado de controle e classes inteiras de
instruções.

\emph{Controle e status.} As instruções de ponto flutuante do RVV leem o modo de arredondamento
\code{frm} (3 bits, cinco modos) e acumulam em \code{fflags} (5 bits), ambos no \code{fcsr}, que precisa
ser salvo, restaurado e renomeado por um núcleo fora de ordem. A aritmética takum define um único modo de
arredondamento e sinaliza resultados excepcionais em banda, com o NaR, de modo que a Zvtakum não lê nem
escreve nenhum CSR além de \code{vl} e \code{vtype}. Em nossas execuções, os núcleos IEEE levantaram flags
em até \FLAGPCT\% das operações sobre elementos (\FLAGWHERE), e o FP8 em \emph{wide} terminou com inexato,
\emph{underflow}, \emph{overflow} e inválido ligados; os núcleos takum não têm esse canal lateral para
manter.

\emph{Instruções que viram operações inteiras.} Como as palavras takum se ordenam como inteiros em
complemento de dois e a negação é a negação inteira, as seis comparações vetoriais de ponto flutuante
(\code{vmfeq}, \code{vmfne}, \code{vmflt}, \code{vmfle}, \code{vmfgt}, \code{vmfge}), mínimo e máximo
(\code{vfmin}, \code{vfmax}), as três instruções de injeção de sinal (\code{vfsgnj}, \code{vfsgnjn},
\code{vfsgnjx}) e \code{vfclass} não precisam de equivalente takum: as comparações inteiras existentes,
\code{vmin}/\code{vmax} e \code{vneg} já dão a resposta certa, com o padrão do NaR ordenado abaixo de
todos os números. O IEEE precisa de instruções dedicadas por causa dos zeros com sinal, dos NaN não
ordenados e da codificação em sinal e magnitude.

\emph{Conversões.} Estreitar é arredondar a palavra (Tabela~\ref{tab:hw}) e alargar é acrescentar bits
zero, o que é exato. As três larguras takum formam, portanto, uma única família de formatos, enquanto
E4M3, E5M2, binary16, bfloat16 e binary32 precisam, cada um, de lógica própria de conversão para cada
outro.

\subsection{Limitações}

O caminho de dados é combinacional e não passou por posicionamento e roteamento; não relatamos
temporização nem uma comparação com uma ALU IEEE de mesma largura. O T16 e o T32 são verificados por
amostragem dirigida e aleatória, não exaustivamente. Todas as matrizes, exceto \emph{memplus}, são pequenas (no máximo 494 linhas) e, nas
reais, $x$ é aleatório; o simulador conta instruções, não ciclos. O ponto flutuante de 8 bits tem a divisão de campos do E4M3 sob as regras do IEEE~754, com infinitos.
Escolhemos assim porque a comparação é contra o padrão IEEE~754, que não define formato de 8 bits, e isso
mantém FP8, FP16 e FP32 como uma única aritmética em três larguras; ele difere do OFP8, cujo E4M3 não tem
infinitos, de modo que parte dos resultados não finitos relatados para o FP8 seria, lá, saturação.

%======================================================================
\section{Conclusão}

Implementamos e verificamos um caminho de dados vetorial takum para o RISC-V e o avaliamos em nível de
instrução. Três propriedades do formato se traduzem diretamente em simplificações de hardware e de ISA: o
cabeçalho fixo de 12 bits dá um único decodificador para todas as larguras; os padrões de bits monótonos
transformam arredondamento, redução de largura, comparação e negação em operações inteiras; e a ausência
de flags e de modos de arredondamento elimina o CSR de ponto flutuante. Numericamente, os takums são a
melhor escolha em 32 bits e, para dados próximos da escala unitária, em 16 bits; em 8 bits, seu valor está
na faixa, não na precisão. Um produto escalar fundido com acumulador de largura dupla se ajusta bem ao
formato, porque os produtos são exatos em largura dupla, e encurta o núcleo de matriz esparsa por vetor.
Como trabalhos futuros, ficam uma implementação com \emph{pipeline} e resultados de temporização, uma
comparação de área em igualdade de condições com unidades IEEE e um conjunto maior de matrizes grandes.

%======================================================================
\begin{thebibliography}{00}
\bibitem{hunhold2025avx} L.~Hunhold, ``Streamlining SIMD ISA extensions with takum arithmetic: A case
study on Intel AVX10.2,'' in \emph{Proc.\ 14th Int.\ Conf.\ Modern Circuits and Systems Technologies
(MOCAST)}, 2025, pp.~1--6.
\bibitem{hunhold2024takum} L.~Hunhold, ``Beating posits at their own game: Takum arithmetic,'' in
\emph{Next Generation Arithmetic (CoNGA 2024)}, LNCS 14666. Springer, 2024, pp.~1--51.
\bibitem{ofp8} P.~Micikevicius \emph{et al.}, ``OCP 8-bit floating point specification (OFP8),'' Open
Compute Project, 2023.
\bibitem{gustafson2017} J.~L. Gustafson and I.~Yonemoto, ``Beating floating point at its own game:
Posit arithmetic,'' \emph{Supercomputing Frontiers and Innovations}, vol.~4, no.~2, pp.~71--86, 2017.
\bibitem{positstd} Posit Working Group, ``Standard for posit arithmetic (2022),'' 2022.
\bibitem{hunhold2024codec} L.~Hunhold, ``Design and implementation of a takum arithmetic hardware codec
in VHDL,'' arXiv:2408.10594, 2024.
\bibitem{rvv} RISC-V International, ``RISC-V `V' vector extension,'' version 1.0, 2021.
\bibitem{hunhold2025solvers} L.~Hunhold and J.~Quinlan, ``Evaluation of bfloat16, posit, and takum
arithmetics in sparse linear solvers,'' in \emph{Proc.\ IEEE Symp.\ Computer Arithmetic (ARITH)}, 2025,
arXiv:2412.20268.
\bibitem{ieee754} \emph{IEEE Standard for Floating-Point Arithmetic}, IEEE Std 754-2019, 2019.
\bibitem{libtakum} L.~Hunhold, ``libtakum,'' versão 2.0.0, 2025. [Online]. Disponível em:
\url{https://github.com/takum-arithmetic/libtakum}
\bibitem{duff1989} I.~S. Duff, R.~G. Grimes, and J.~G. Lewis, ``Sparse matrix test problems,''
\emph{ACM Trans.\ Math.\ Softw.}, vol.~15, no.~1, pp.~1--14, 1989.
\bibitem{boisvert1997} R.~F. Boisvert, R.~Pozo, K.~Remington, R.~F. Barrett, and J.~J. Dongarra, ``Matrix
Market: A web resource for test matrix collections,'' in \emph{Quality of Numerical Software}.
Chapman \& Hall, 1997, pp.~125--137.
\end{thebibliography}

\end{document}
